%% This document created by Scientific Word (R) Version 3.5

\documentclass[aps,preprint]{revtex4}%
\usepackage{amsmath}
\usepackage{graphicx}
\usepackage{amsfonts}
\usepackage{amssymb}%
\setcounter{MaxMatrixCols}{30}
%TCIDATA{OutputFilter=latex2.dll}
%TCIDATA{Version=4.00.0.2321}
%TCIDATA{CSTFile=revtex4.cst}
%TCIDATA{Created=Tuesday, April 30, 2002 08:45:12}
%TCIDATA{LastRevised=Friday, May 03, 2002 17:24:12}
%TCIDATA{<META NAME="GraphicsSave" CONTENT="32">}
%TCIDATA{<META NAME="DocumentShell" CONTENT="Articles\SW\REVTeX 4">}
\newtheorem{theorem}{Theorem}
\newtheorem{acknowledgement}[theorem]{Acknowledgement}
\newtheorem{algorithm}[theorem]{Algorithm}
\newtheorem{axiom}[theorem]{Axiom}
\newtheorem{claim}[theorem]{Claim}
\newtheorem{conclusion}[theorem]{Conclusion}
\newtheorem{condition}[theorem]{Condition}
\newtheorem{conjecture}[theorem]{Conjecture}
\newtheorem{corollary}[theorem]{Corollary}
\newtheorem{criterion}[theorem]{Criterion}
\newtheorem{definition}[theorem]{Definition}
\newtheorem{example}[theorem]{Example}
\newtheorem{exercise}[theorem]{Exercise}
\newtheorem{lemma}[theorem]{Lemma}
\newtheorem{notation}[theorem]{Notation}
\newtheorem{problem}[theorem]{Problem}
\newtheorem{proposition}[theorem]{Proposition}
\newtheorem{remark}[theorem]{Remark}
\newtheorem{solution}[theorem]{Solution}
\newtheorem{summary}[theorem]{Summary}
\newenvironment{proof}[1][Proof]{\noindent\textbf{#1.} }{\ \rule{0.5em}{0.5em}}
\begin{document}
\preprint{HEP/123-qed}
\title[Short title for running header]{Long title that appears on the title page}
\author{First Author}
\affiliation{My Institution}
\author{Second Author}
\affiliation{My Institution}
\author{Third Author}
\affiliation{Other Institution}
\keywords{one two three}
\pacs{PACS number}

\begin{abstract}
Shell document for REV\TeX {} 4.

\end{abstract}
\volumeyear{year}
\volumenumber{number}
\issuenumber{number}
\eid{identifier}
\date[Date text]{date}
\received{date}

\revised{date}

\accepted{date}

\published{date}

\startpage{101}
\endpage{102}
\maketitle
\tableofcontents


\section{Properties of Diagonal Operators}

\bigskip As
\begin{equation}
\left\langle \mathsf{x}^{N}|\mathsf{y}^{N}\right\rangle =\delta\left(
\mathsf{x}^{N}-\mathsf{y}^{N}\right)  \label{on}%
\end{equation}
one can deduce that
\begin{equation}
e^{\breve{\varsigma}_{N}}=e^{\int\varsigma\left(  \mathsf{z}^{N}\right)
\left|  \mathsf{z}^{N}\right\rangle \left\langle \mathsf{z}^{N}\right|
d\mathsf{z}^{N}}=\int e^{\varsigma\left(  \mathsf{z}^{N}\right)  }\left|
\mathsf{z}^{N}\right\rangle \left\langle \mathsf{z}^{N}\right|  d\mathsf{z}%
^{N}%
\end{equation}
We can decompose $\varsigma\left(  \mathsf{z}^{N}\right)  $ into irreducible
(pure) parts as
\begin{equation}
\varsigma\left(  \mathsf{z}^{N}\right)  =\sum_{1\leq k\leq N}\varsigma
_{k}\left(  \mathsf{z}^{N}\right)
\end{equation}
where
\begin{equation}
\varsigma_{k}\left(  \mathsf{z}^{N}\right)  =\sum_{1\leq i_{1}<i_{2}%
\cdots<i_{k}\leq N}\varsigma_{k}\left(  \mathsf{z}_{i_{1}}\mathsf{z}_{i_{2}%
}\cdots\mathsf{z}_{i_{k}}\right)
\end{equation}
$\lceil\varsigma_{k}\left(  \mathsf{z}_{i_{1}}\mathsf{z}_{i_{2}}%
\cdots\mathsf{z}_{i_{k}}\right)  $ are all symmetric functions$\rfloor$ so
that
\begin{align}
\int\varsigma\left(  \mathsf{z}^{N}\right)  \left|  \mathsf{z}^{N}%
\right\rangle \left\langle \mathsf{z}^{N}\right|  d\mathsf{z}^{N}  &
=\int\sum_{1\leq k\leq N}\varsigma_{k}\left(  \mathsf{z}^{N}\right)  \left|
\mathsf{z}^{N}\right\rangle \left\langle \mathsf{z}^{N}\right|  d\mathsf{z}%
^{N}\nonumber\\
&  =\int\sum_{1\leq k\leq N}\sum_{1\leq i_{1}<i_{2}\cdots<i_{k}\leq
N}\varsigma_{k}\left(  \mathsf{z}_{i_{1}}\mathsf{z}_{i_{2}}\cdots
\mathsf{z}_{i_{k}}\right)  \left|  \mathsf{z}^{N}\right\rangle \left\langle
\mathsf{z}^{N}\right|  d\mathsf{z}^{N}\nonumber\\
&  =\sum_{1\leq k\leq N}\sum_{1\leq i_{1}<i_{2}\cdots<i_{k}\leq N}%
\int\varsigma_{k}\left(  \mathsf{z}_{i_{1}}\mathsf{z}_{i_{2}}\cdots
\mathsf{z}_{i_{k}}\right)  \left|  \mathsf{z}^{N}\right\rangle \left\langle
\mathsf{z}^{N}\right|  d\mathsf{z}^{N}%
\end{align}
$\lceil$If the orthonormalization condition \ eq. $\left(  \ref{on}\right)  $
is not satisfied e.g. when the overcomplete basis is formed by proper vectors
we have
\begin{equation}
\left\langle \mathsf{x}^{N}|\mathsf{y}^{N}\right\rangle =K\left(
\mathsf{x}^{N},\mathsf{y}^{N}\right)
\end{equation}
and
\begin{align}
e^{\breve{\varsigma}_{N}}  &  =e^{\int\varsigma\left(  \mathsf{z}^{N}\right)
\left|  \mathsf{z}^{N}\right\rangle \left\langle \mathsf{z}^{N}\right|
d\mathsf{z}^{N}}=1+\int\varsigma\left(  \mathsf{z}^{N}\right)  \left|
\mathsf{z}^{N}\right\rangle \left\langle \mathsf{z}^{N}\right|  d\mathsf{z}%
^{N}\nonumber\\
&  \qquad\qquad\qquad\qquad+\frac{1}{2}\int\varsigma\left(  \mathsf{z}_{1}%
^{N}\right)  K\left(  \mathsf{z}_{1}^{N},\mathsf{z}_{2}^{N}\right)
\varsigma\left(  \mathsf{z}_{2}^{N}\right)  \left|  \mathsf{z}_{1}%
^{N}\right\rangle \left\langle \mathsf{z}_{2}^{N}\right|  d\mathsf{z}_{1}%
^{N}d\mathsf{z}_{2}^{N}+\cdots
\end{align}
$\rfloor$

The operator
\begin{equation}
\breve{\varsigma}_{N}=\int\varsigma\left(  \mathsf{z}^{N}\right)  \left\vert
\mathsf{z}^{N}\right\rangle \left\langle \mathsf{z}^{N}\right\vert
d\mathsf{z}^{N}%
\end{equation}
is the restriction of the second quantized operator%
\begin{equation}
\breve{\varsigma}=\int\varsigma\left(  \mathsf{z}^{N}\right)  \Phi\left(
\mathsf{z}^{N}\right)  ^{\dagger}\Phi\left(  \mathsf{z}^{N}\right)
d\mathsf{z}^{N}%
\end{equation}
to $\mathcal{H}^{N}$, where
\begin{equation}
\Phi\left(  \mathsf{z}^{N}\right)  =\Phi\left(  \mathsf{z}_{N}\right)
\cdots\Phi\left(  \mathsf{z}_{1}\right)
\end{equation}
and
\begin{equation}
\Phi\left(  \mathsf{z}\right)  ^{\dagger}\left\vert \phi\right\rangle
=\left\vert \mathsf{z}\right\rangle
\end{equation}
First quantized operators can always be thought of as projections of number
preserving second quantized operators i.e.
\begin{align}
P_{N}\mathfrak{A}P_{N} &  =P_{N}\int A\left(  \mathsf{x}^{N},\mathsf{y}%
^{N}\right)  \Phi\left(  \mathsf{x}^{N}\right)  ^{\dagger}\Phi\left(
\mathsf{y}^{N}\right)  d\mathsf{x}^{N}d\mathsf{y}^{N}P_{N}\nonumber\\
&  =\int A\left(  \mathsf{x}^{N},\mathsf{y}^{N}\right)  \left\vert
\mathsf{x}^{N}\right\rangle \left\langle \mathsf{y}^{N}\right\vert
d\mathsf{x}^{N}d\mathsf{y}^{N}%
\end{align}
so that
\begin{equation}
\mathfrak{A=}\sum_{N=0}^{\infty}P_{N}\mathfrak{A}P_{N}%
\end{equation}
The question one can then ask is whether
\begin{equation}
F\left(  P_{N}\mathfrak{A}P_{N}\right)  =P_{N}F\left(  \mathfrak{A}\right)
P_{N}%
\end{equation}
in particular for the exponential function. The answer is obtained by
considering the following
\begin{align}
e^{\sum_{N=0}^{\infty}P_{N}\mathfrak{A}P_{N}} &  =e^{\mathfrak{A}}\nonumber\\
&  =\sum_{M=0}^{\infty}P_{M}+\sum_{M=0}^{\infty}P_{M}\mathfrak{A}P_{M}%
+\frac{1}{2}\sum_{M,M^{\prime}=0}^{\infty}P_{M}\mathfrak{A}P_{M}P_{M^{\prime}%
}\mathfrak{A}P_{M^{\prime}}+\cdots\nonumber\\
&  =\sum_{M=0}^{\infty}P_{M}+\sum_{M=0}^{\infty}P_{M}\mathfrak{A}P_{M}%
+\frac{1}{2}\sum_{M=0}^{\infty}P_{M}\mathfrak{A}^{2}P_{M}+\cdots
\end{align}
Hence
\begin{align}
P_{N}e^{\mathfrak{A}}P_{N} &  =P_{N}+P_{N}\mathfrak{A}P_{N}+\frac{1}{2}%
\sum_{N=0}^{\infty}P_{N}\mathfrak{A}^{2}P_{N}+\cdots\nonumber\\
&  =P_{N}e^{P_{N}\mathfrak{A}P_{N}}%
\end{align}
Thus
\begin{align}
P_{N}e^{\mathfrak{A}}P_{N} &  =P_{N}e^{\int\varsigma\left(  \mathsf{z}%
^{N}\right)  P_{N}\Phi\left(  \mathsf{z}^{N}\right)  ^{\dagger}\Phi\left(
\mathsf{z}^{N}\right)  P_{N}d\mathsf{z}^{N}}\nonumber\\
&  =P_{N}e^{\int\varsigma\left(  \mathsf{z}^{N}\right)  \left\vert
\mathsf{z}^{N}\right\rangle \left\langle \mathsf{z}^{N}\right\vert
d\mathsf{z}^{N}}%
\end{align}
When $\mathsf{z}_{i}\neq\mathsf{z}_{j},\quad1\leq i<j\leq N$ the product of
field operators can always be rearranged as
\begin{align}
\Phi\left(  \mathsf{z}^{N}\right)  ^{\dagger}\Phi\left(  \mathsf{z}%
^{N}\right)   &  =\Phi\left(  \mathsf{z}_{1}\right)  ^{\dagger}\cdots
\Phi\left(  \mathsf{z}_{N}\right)  ^{\dagger}\Phi\left(  \mathsf{z}%
_{N}\right)  \cdots\Phi\left(  \mathsf{z}_{1}\right)  \nonumber\\
&  =\Phi\left(  \mathsf{z}_{1}\right)  ^{\dagger}\Phi\left(  \mathsf{z}%
_{1}\right)  \cdots\Phi\left(  \mathsf{z}_{N}\right)  ^{\dagger}\Phi\left(
\mathsf{z}_{N}\right)
\end{align}
and%
\begin{align}
\breve{\varsigma} &  =\int\varsigma\left(  \mathsf{z}^{N}\right)  \prod_{1\leq
i<j\leq N}\delta^{C}\left(  \mathsf{z}_{i}-\mathsf{z}_{j}\right)  \Phi\left(
\mathsf{z}_{1}\right)  ^{\dagger}\Phi\left(  \mathsf{z}_{1}\right)  \cdots
\Phi\left(  \mathsf{z}_{N}\right)  ^{\dagger}\Phi\left(  \mathsf{z}%
_{N}\right)  d\mathsf{z}^{N}\nonumber\\
&  =\int\sum_{k=0}^{N}\varsigma_{k}\left(  \mathsf{z}^{N}\right)  \prod_{1\leq
i<j\leq N}\delta^{C}\left(  \mathsf{z}_{i}-\mathsf{z}_{j}\right)  \Phi\left(
\mathsf{z}_{1}\right)  ^{\dagger}\Phi\left(  \mathsf{z}_{1}\right)  \cdots
\Phi\left(  \mathsf{z}_{N}\right)  ^{\dagger}\Phi\left(  \mathsf{z}%
_{N}\right)  d\mathsf{z}^{N}%
\end{align}
where $\delta^{C}\left(  \mathsf{z}_{i}-\mathsf{z}_{j}\right)  =1-\delta
\left(  \mathsf{z}_{i}-\mathsf{z}_{j}\right)  $ and
\begin{equation}
\varsigma_{k}\left(  \mathsf{z}^{N}\right)  =\sum_{1\leq i_{1}<i_{2}%
\cdots<i_{k}\leq N}\varsigma_{k}\left(  \mathsf{z}_{i_{1}}\mathsf{z}_{i_{2}%
}\cdots\mathsf{z}_{i_{k}}\right)
\end{equation}
$\lceil$This gives for example
\begin{equation}
\varsigma_{3}\left(  \mathsf{z}^{N}\right)  =\sum_{1\leq i_{1}<i_{2}<i_{3}\leq
N}\varsigma_{3}\left(  \mathsf{z}_{i_{1}}\mathsf{z}_{i_{2}}\mathsf{z}_{i_{3}%
}\right)
\end{equation}
and
\begin{align}
\breve{\varsigma}_{3} &  =\int\varsigma_{3}\left(  \mathsf{z}^{N}\right)
\prod_{1\leq i<j\leq N}\delta^{C}\left(  \mathsf{z}_{i}-\mathsf{z}_{j}\right)
\Phi\left(  \mathsf{z}_{1}\right)  ^{\dagger}\Phi\left(  \mathsf{z}%
_{1}\right)  \cdots\Phi\left(  \mathsf{z}_{N}\right)  ^{\dagger}\Phi\left(
\mathsf{z}_{N}\right)  d\mathsf{z}^{N}\nonumber\\
&  =\int\sum_{1\leq i_{1}<i_{2}<i_{3}\leq N}\varsigma_{3}\left(
\mathsf{z}_{i_{1}}\mathsf{z}_{i_{2}}\mathsf{z}_{i_{3}}\right)  \prod_{1\leq
i<j\leq N}\delta^{C}\left(  \mathsf{z}_{i}-\mathsf{z}_{j}\right)  \Phi\left(
\mathsf{z}_{1}\right)  ^{\dagger}\Phi\left(  \mathsf{z}_{1}\right)  \cdots
\Phi\left(  \mathsf{z}_{N}\right)  ^{\dagger}\Phi\left(  \mathsf{z}%
_{N}\right)  d\mathsf{z}^{N}\nonumber\\
&  =\int\sum_{1\leq i_{1}<i_{2}<i_{3}\leq N}\varsigma_{3}\left(
\mathsf{z}_{i_{1}}\mathsf{z}_{i_{2}}\mathsf{z}_{i_{3}}\right)  \prod_{1\leq
i<j\leq N}\delta^{C}\left(  \mathsf{z}_{i}-\mathsf{z}_{j}\right)  \Phi\left(
\mathsf{z}_{i_{1}}\right)  ^{\dagger}\Phi\left(  \mathsf{z}_{i_{1}}\right)
\Phi\left(  \mathsf{z}_{i_{2}}\right)  ^{\dagger}\Phi\left(  \mathsf{z}%
_{i_{2}}\right)  \Phi\left(  \mathsf{z}_{i_{3}}\right)  ^{\dagger}\Phi\left(
\mathsf{z}_{i_{3}}\right)  \nonumber\\
&  \qquad\qquad\qquad\qquad\qquad\qquad\qquad\qquad\qquad\qquad\qquad
\times\Phi\left(  \mathsf{z}_{i_{4}}\right)  ^{\dagger}\Phi\left(
\mathsf{z}_{i_{4}}\right)  \cdots\Phi\left(  \mathsf{z}_{i_{N}}\right)
^{\dagger}\Phi\left(  \mathsf{z}_{i_{N}}\right)  d\mathsf{z}^{N}%
\end{align}
$\rfloor$

Thus we have the following
\begin{align}
e^{\breve{\varsigma}_{N}} &  =e^{\int\varsigma\left(  \mathsf{z}^{N}\right)
\left\vert \mathsf{z}^{N}\right\rangle \left\langle \mathsf{z}^{N}\right\vert
d\mathsf{z}^{N}}\nonumber\\
&  =P_{N}e^{\int\varsigma\left(  \mathsf{z}^{N}\right)  \prod_{1\leq i<j\leq
N}\delta^{C}\left(  \mathsf{z}_{i}-\mathsf{z}_{j}\right)  P_{N}\Phi\left(
\mathsf{z}_{1}\right)  ^{\dagger}\Phi\left(  \mathsf{z}_{1}\right)  \cdots
\Phi\left(  \mathsf{z}_{N}\right)  ^{\dagger}\Phi\left(  \mathsf{z}%
_{N}\right)  P_{N}d\mathsf{z}^{N}}\nonumber\\
&  =\int e^{\varsigma\left(  \mathsf{z}^{N}\right)  }\left\vert \mathsf{z}%
^{N}\right\rangle \left\langle \mathsf{z}^{N}\right\vert d\mathsf{z}%
^{N}\nonumber\\
&  =\int e^{\varsigma\left(  \mathsf{z}^{N}\right)  }\prod_{1\leq i<j\leq
N}\delta^{C}\left(  \mathsf{z}_{i}-\mathsf{z}_{j}\right)  P_{N}\Phi\left(
\mathsf{z}_{1}\right)  ^{\dagger}\Phi\left(  \mathsf{z}_{1}\right)  \cdots
\Phi\left(  \mathsf{z}_{N}\right)  ^{\dagger}\Phi\left(  \mathsf{z}%
_{N}\right)  P_{N}d\mathsf{z}^{N}\nonumber\\
&  =\int\left\{  1+\varsigma\left(  \mathsf{z}^{N}\right)  +\varsigma\left(
\mathsf{z}^{N}\right)  ^{2}+\cdots\right\}  \left\vert \mathsf{z}%
^{N}\right\rangle \left\langle \mathsf{z}^{N}\right\vert d\mathsf{z}%
^{N}\nonumber\\
&  =\int\left\{  1+\sum_{k=0}^{N}\sum_{1\leq i_{1}<i_{2}\cdots<i_{k}\leq
N}\varsigma_{k}\left(  \mathsf{z}_{i_{1}}\mathsf{z}_{i_{2}}\cdots
\mathsf{z}_{i_{k}}\right)  +\varsigma\left(  \mathsf{z}^{N}\right)
^{2}\right.  \nonumber\\
&  \qquad\qquad\left.  +\left\{  \sum_{k=0}^{N}\sum_{1\leq i_{1}<i_{2}%
\cdots<i_{k}\leq N}\varsigma_{k}\left(  \mathsf{z}_{i_{1}}\mathsf{z}_{i_{2}%
}\cdots\mathsf{z}_{i_{k}}\right)  \right\}  ^{2}+\cdots\right\}  \left\vert
\mathsf{z}^{N}\right\rangle \left\langle \mathsf{z}^{N}\right\vert
d\mathsf{z}^{N}%
\end{align}


For example we can consider the case of $N=2$
\begin{align}
e^{\breve{\varsigma}_{N}}  & =\int\left\{  1+\alpha+\sum_{1\leq i_{1}%
<i_{2}<\leq N}\left\{  \varsigma_{1}\left(  \mathsf{z}_{i_{1}}\right)
+\varsigma_{1}\left(  \mathsf{z}_{i_{2}}\right)  +\varsigma_{2}\left(
\mathsf{z}_{i_{1}}\mathsf{z}_{i_{2}}\right)  \right\}  \right.  \nonumber\\
& \left.  +\frac{1}{2}\left\{  \alpha+\sum_{1\leq i_{1}<i_{2}<\leq N}\left\{
\varsigma_{1}\left(  \mathsf{z}_{i_{1}}\right)  +\varsigma_{1}\left(
\mathsf{z}_{i_{2}}\right)  +\varsigma_{2}\left(  \mathsf{z}_{i_{1}}%
\mathsf{z}_{i_{2}}\right)  \right\}  \right\}  ^{2}+\cdots\right\}  \left\vert
\mathsf{z}^{N}\right\rangle \left\langle \mathsf{z}^{N}\right\vert
d\mathsf{z}^{N}%
\end{align}



\end{document}